\section{Twisted pseudo-identity}\label{sec:pseudoidentity}

The ordinary and extraordinary diagonal kernels play different roles in the
argument.  The former represents categorical duality, as established in
Section~\ref{sec:foundations}; the latter must itself define an equivalence.
We prove this by combining parabolic adjunctions with a compactification of
the diagonal.  The essential additional point at an irrational twist is that
central inertia removes every boundary stratum with nonzero defect.

Let $H$ be a connected reductive group, let $V$ be a self-dual representation
whose trace form $Q$ is positive definite on the real cocharacter space, and
put $L_H=\det R\Gamma(X,V_E)$.  Restrictions of $V$ will define the corresponding
lines for Levi subgroups.  Write
\[
 \mathcal C_H=D(L_H^s),\qquad \mathcal C'_H=D(L_H^{-s}),\qquad s\notin\mathbb Q.
\]
The duality of Section~\ref{sec:foundations} identifies
$(\mathcal C'_H)^\vee$ with $\mathcal C_H$.  The line ratio on the diagonal is
canonically trivial.  Consequently the kernel $\Delta_!k_{Y_H}$ defines a
continuous endofunctor $F_H$ of $\mathcal C_H$; its existence as a kernel uses
that the diagonal is schematic and that $k_{Y_H}$ is locally holonomic.

\begin{theorem}\label{pseudo:equivalence}
For every $H,V,s$ as above, the functor $F_H$ is an equivalence.  The assertion
holds on all components, including those on which the twisted category is
zero, and also with $s$ replaced by $-s$.
\end{theorem}

We first establish parabolic adjunctions and their compatibility with
$F_H$.  The compactified diagonal will then show that the cone of the
natural map $F_H\to\operatorname{Id}[-2\dim Y_H]$ has a finite
filtration whose quotients factor through proper Levi categories.  Induction handles objects obtained by
Eisenstein series, while $F_H$ acts by a shift on objects whose proper
constant terms vanish.  These two cases generate the category and will
prove the equivalence.

The geometric constructions are those of geometric Eisenstein series and second
adjointness \cite{BravermanGaitsgory,DrinfeldGaitsgoryCT}.  We include the twist
and coherence arguments, since an adjunction only on holonomic objects would
not suffice here.

\subsection{Parabolic correspondences and their compactifications}

For the moment suppose that $H$ is semisimple.  Fix a parabolic $P$, its Levi
quotient $M$, an opposite $P^-$, and a torsion degree $u\in\pi_1(M)$.  Set
$Y=Y_H$ and use the correspondence
\[
 Y_M^u\xleftarrow{\ q\ }Y_P^u\xrightarrow{\ p\ }Y.
\]
The weight filtration of $V$ associated to a cocharacter defining $P$ gives
an identification $p^*L_H\simeq q^*L_M$.  Normalize it to be the identity on
split $M$-reductions.  Use the same normalization for $P^-$.  Define
\[
 C_P=q_*p^!,\qquad E_P=p_*q^!,\qquad
 e=\dim(q)=(g(X)-1)\dim N_P.
\]
The last equality uses that $u$ is torsion; it gives the same value for $P^-$.
Primes on these functors will mean that they are taken at the opposite twist.

Here $p$ is schematic, separated, and relatively quasi-compact, whereas $q$
is smooth, relatively quasi-compact, and safe.  A positive grading of $N_P$
presents $q$, locally over the Levi stack, as successive fibrations of vector
stacks.  A vector-stack step has the form $[V^1/V^0]$, with $V^i$ vector bundles
coming from a two-term complex computing cohomology on $X$.  There are no
obstructions in degree two.  This description has three useful consequences.
First, the direct images defining $C_P,E_P$ are continuous.  Second, $q_*$
preserves local regular holonomicity.  Third, the partial holonomic operations
$q_!$ and $q^*$ have their actual adjunction property against arbitrary
D-modules.  For the last two assertions present each vector stack by its
vector bundle: descent along the vector group is fully faithful by its de
Rham acyclicity, and both assertions reduce to the schematic formulas.
More explicitly, for $a:V^1\to[V^1/V^0]$, the star image is
computed using $a^*$ followed by the schematic star image; the holonomic
shriek image uses $a^!$ followed by the schematic shriek image.  The
identities $a_*a^*\simeq\operatorname{Id}$ and
$a_!a^!\simeq\operatorname{Id}$ follow from the corresponding ordinary
and compactly supported cohomology of the vector-group fibers.  This
retains the shifts even when the vector stack has negative dimension.
These reductions are compatible with smooth descent and with adding a
spectator category.

Kernel duality gives
\[
 (C'_P)^\vee=E_P.
\]
Indeed both kernels are obtained from the diagonal unit by !-base change
followed by the indicated safe direct image.  Every calculation below may
be made on quasi-compact open exhaustions; the boundedness result of
Section~\ref{sec:foundations} identifies the twisted categories of the bundle
stacks with those on fixed quasi-compact opens.  Thus this calculation uses
ordinary categorical duals and does not introduce a co-category.

We will use the fine Drinfeld compactification
\[
 Y_P^u\xrightarrow{\ j\ }\widetilde Y_P^u,
 \qquad \bar q:\widetilde Y_P^u\longrightarrow Y_M^u,
 \qquad \bar p:\widetilde Y_P^u\longrightarrow Y.
\]
Its points are generalized reductions with their Levi datum retained, and
$\bar p$ is schematic and proper.  If the derived group is simply connected,
this is the compactification by generically nondegenerate sections of the
affine closure of $H/N_P$.  The boundary is stratified by positive defects.
On each stratum the saturated reduction is an actual $P$-bundle; its Levi
bundle is related to the originally specified one by a positive
modification.  The map from the stratum to this Levi Hecke datum is the
base change of the ordinary map from parabolic bundles to Levi bundles.
The stratification is locally finite on quasi-compact substacks.  These are
the geometric compactification and stratification results of
\cite{BravermanGaitsgory}; see also
\cite[Sections~3.2--3.3]{SchiederVinberg}.

For completeness, the same statements hold for the actual global form of
$H$.  Choose a central extension $H_1\to H$ with torus kernel $D$ and simply
connected derived group, and form the compactification upstairs.  Central
$D$-bundles act on all its data: each Pl\"ucker block transforms through its
central character.  Quotient by this action.  Lifts of an $H$-bundle exist
smoothly locally on the base.  On geometric fibers this follows from the
vanishing of the Brauer group of the smooth projective curve; infinitesimal
lifting follows from the vanishing of degree-two coherent cohomology.
Thus $Y_{H_1}\to Y_H$ and the corresponding parabolic and Levi maps are
torsors for $Y_D$ in this sense.

Fix a lift of the ambient degree and the original Levi degree downstairs.
There is a unique compatible degree for its lift.  To check this, use the
exact sequence of maximal tori: two possible Levi lifting degrees differ
by a cocharacter of $D$, and this cocharacter injects into $\pi_1(H_1)$.
The injectivity follows because $D$ meets the derived group in a finite
subgroup.  Compatibility of ambient degrees therefore forces the difference
to vanish.  Conversely the same sequence provides any compatible lift.
Properness, schematicity, and the stratum description consequently descend.
The descended Hecke base is the simultaneous central quotient of the
upstairs Hecke base; its unipotent-extension fibers are unchanged.  A
nonzero defect changes the Levi degree nontrivially over $\mathbb Q$: the
simple coroots omitted from $M$ remain linearly independent in the quotient
of the rational cocharacter space by the Levi coroot span.  This rational
statement, rather than a possibly enlarged integral defect monoid, is what
we shall use.

\subsection{Boundary smoothness and cleanness}

On $\widetilde Y_P^u$ take the twist
\[
 (\bar p^*L_H\otimes\bar q^*L_M^{-1})^s
\]
and the kernel $J=j_!\omega_{Y_P^u}$ defined using its splitting over the open.
We next prove two properties: $J$ is clean, and it has no nonzero
characteristic covector pulled back from the Levi base.  The second property
allows us to tensor $J$ with arbitrary coherent Levi coefficients.

Here is the elementary analytic observation that will be used twice.  Let
$U\subset Z$ be an open subset of a finite-type complex scheme, let
$\mathcal L$ be a local system on $U$, and let $j_!\mathcal L$ be its extension
by zero.  A derivation of $\mathcal O_Z$ preserving the ideal of $Z\setminus U$
has local analytic flows preserving the pair $(Z,U)$.  To see this even when
$Z$ is singular, embed it locally as a closed subscheme of a smooth scheme.
Lift the derivation to the ambient local coordinates.  Its preservation of
the ideals of $Z$ and of the boundary makes the ambient flow preserve their
analytic zero loci.  Along a sufficiently small time disk, flow transport
identifies the two pullbacks of $\mathcal L$, and hence those of
$j_!\mathcal L$.  These identifications commute with extraordinary pullback
along a compatible map of pairs.  On a smooth test scheme, a family of such
fields spanning the tangent space therefore makes the pulled-back complex
a local system: successive flows give local coordinates.  Likewise every
characteristic covector of the closed direct image annihilates each such
field.  This last assertion follows from the equality of characteristic
variety and microsupport for regular holonomic D-modules
\cite[Theorem~11.3.3]{KashiwaraSchapira} and the local product description
along the flow.

\begin{lemma}\label{pseudo:fiber-descent}
Let $i_\theta:S_\theta\hookrightarrow\widetilde Y_P^u$ be a defect stratum,
and let $f_\theta:S_\theta\to B_\theta$ be its Levi Hecke-data projection.
Its twist is pulled back from $B_\theta$, and there is a locally holonomic
complex $K_\theta$ on that base such that
\begin{equation}\label{eq:M2}
 i_\theta^!J\simeq f_\theta^!K_\theta.
\end{equation}
\end{lemma}

\begin{proof}
The determinant filtration for the saturated reduction identifies the line
ratio with the ratio of the two Levi determinant lines.  This proves the
assertion about twists.  We first show smoothness on each geometric
!-fiber of $f_\theta$.

Choose a point on the stratum, and a point $z\in X$ outside its defect divisor.
Shrink the parameter chart so that $z$ remains outside every defect.  On the
formal disk at $z$ the generalized reduction is actual.  Choose a frame of
this reduction, locally on a smooth chart; for any finite collection of
Laurent principal parts only a finite jet of this choice is needed.  Frames
of that jet lift smoothly, and compatible higher jets exist by formal
smoothness.  In the central-extension construction take the same charts
upstairs before descent.

A Laurent element $\eta\in\mathfrak n_P((t_z))$ modifies the framed bundle
infinitesimally by $1+\epsilon\eta$.  Gluing on the punctured disk produces a
first-order deformation of the generalized reduction, unchanged away from
$z$.  Since the data are actual near $z$, this gluing preserves their defining
relations.  A smooth scheme presentation $T'\to\widetilde Y_P^u$ admits a
local lift of this infinitesimal deformation to a derivation of $T'$; any
chosen value differing by a relative tangent vector can be prescribed.
The derivation preserves the boundary ideal.  More explicitly, the defect
incidence is unchanged away from $z$ and is empty near $z$.  Its scheme image
in the parameter space is therefore unchanged by the automorphism over
$\mathbb C[\epsilon]/(\epsilon^2)$ associated to the derivation.  Formation
of this image commutes with these flat pullbacks.  This proves ideal
preservation, not just preservation of boundary points.

Fix the Levi Hecke datum and pull $T'$ back to the stratum fiber, obtaining
$W$.  The map from $W$ to that fiber is smooth, being a base change of the
smooth presentation $T'\to\widetilde Y_P^u$.  The fiber itself is a smooth
tower of vector stacks over a point.  Thus $W$ is smooth, after a scheme
presentation when necessary.
The fields just constructed lift compatibly to $W$, because a modification
by $N_P$ keeps the Levi modification fixed.  Their values, together with
relative presentation directions, span its tangent space.  Indeed the
non-presentation directions are deformations of a parabolic bundle with
fixed Levi, hence are the image of
$H^1(X,(\mathfrak n_P)_{E_P})$.  Principal parts at $z$ surject onto this
space: for sufficiently large $n$, the exact sequence for
$(\mathfrak n_P)_{E_P}\subset(\mathfrak n_P)_{E_P}(nz)$ and Serre vanishing
give the surjection.  Its connecting map is precisely the gluing deformation
above.  The tangent-complex sequence supplies the remaining presentation
directions, including the quotient by infinitesimal automorphisms.

Trivialize the line-frame torsors locally analytically.  The kernel $J$ is
regular holonomic, and on the open it corresponds to a local system.
The flow observation now shows that its !-pullback to $W$ is a local system.
Changing a frame only tensors by a smooth Kummer local system, so does not
alter this conclusion.

It remains to turn fiberwise smoothness into descent.  Each step of
$f_\theta$ is a vector-stack fibration.  On a vector stack over a point, a
smooth regular holonomic complex is constant, since both the vector space
and the acting vector group are contractible.  For one fibration apply the
counit $f^*f_*\to\operatorname{Id}$ and test it on geometric !-fibers by
!-base change.  It is an isomorphism on each fiber and therefore globally.
Repeat for the tower and use smooth descent.  Since $f^*$ differs from $f^!$
by the smooth dimension shift, this gives \eqref{eq:M2}.
\end{proof}

\begin{lemma}\label{pseudo:clean-noncharacteristic}
The kernel $J$ is clean: $j_!\omega\simeq j_*\omega$ in the specified twist.
Moreover, on a smooth Levi chart $S\to Y_M^u$, choose a scheme presentation
of $\widetilde Y_P^u\times_{Y_M^u}S$ and a local closed embedding into a
smooth scheme over $S$.  The characteristic variety of the closed direct
image of $J$ contains no nonzero covector pulled back from $S$.
\end{lemma}

\begin{proof}
For $\theta\ne0$, simultaneous central inertia on the two Levi bundles of
$B_\theta$ acts on the determinant ratio with weight
$Q(z,u_{\rm sat}-u)$.  The degree difference is rationally nonzero, so
positivity of $Q$ provides an integral central cocharacter $z$ for which
this integer is nonzero.  At exponent $s$ its Kummer parameter is not
integral.  The inertia vanishing test of Section~\ref{sec:foundations}
therefore kills the entire twisted category on this base, including any
additional spectator.  Under a central extension one may first multiply
and lift $z$; the removed central directions act trivially on the ratio,
so the test descends.  Lemma~\ref{pseudo:fiber-descent} gives
$i_\theta^!J=0$.  Localization and the locally finite stratification prove
cleanness.  The same reasoning applies to the opposite twist and hence to
the Verdier-dual kernel.

For the characteristic assertion repeat the modification argument at $z$
with elements of the Levi loop Lie algebra.  Prescribe their lifts first in
$S$, then in the presentation of the compactification.  These derivations
again preserve the open and its boundary.  Principal parts in the adjoint
Levi bundle span its $H^1$, and the relative chart directions complete their
images to a spanning set of $T_sS$.  Extend the derivations in the chosen
smooth ambient embedding.  The flow observation shows that every
characteristic covector annihilates them.  A horizontal covector is pulled
back from $T_s^*S$; annihilating fields whose projections span $T_sS$ forces
it to be zero.  This proves the assertion.
\end{proof}

\begin{proposition}\label{pseudo:first-adjunction}
On degree $u$, the continuous functor $L_P=E_P[-2e]$ is a left adjoint of
$C_P$:
\begin{equation}\label{eq:M1}
 L_P=E_P[-2e]\ \dashv\ C_P.
\end{equation}
It computes the partially defined expression $p_!q^*$ on every object.
\end{proposition}

\begin{proof}
First let $E\in\mathcal C_M^u$ be compact, hence locally coherent by
Section~\ref{sec:foundations}, and form
\[
 J_E=J\otimes^!\bar q^!E.
\]
In the charts of Lemma~\ref{pseudo:clean-noncharacteristic}, the second
factor has horizontal characteristic covectors.  The absence of nonzero
horizontal covectors for the first factor is exactly the condition making
the diagonal noncharacteristic for their exterior product.  The
noncharacteristic inverse-image theorem, including compatibility with
duality \cite[Theorems~2.4.6 and~2.7.1(ii)]{HottaTakeuchiTanisaki}, proves that $J_E$ is
locally coherent and that Verdier duality commutes with this tensor product,
up to the smooth dimension shift.  The same calculation applies after
closed direct image from the possibly singular presentation, by the
projection formula.

Both $J_E$ and its Verdier dual have zero !-restriction to the boundary,
since this holds for $J$ and its dual.  Hence $J_E$ is simultaneously the
star and the extraordinary extension of its restriction to $Y_P^u$.
Here extraordinary extension has its full partial-adjoint meaning:
there are no maps from $J_E$ to an arbitrary object pushed from the
boundary.  We verify this for arbitrary targets using chartwise compactness
and descent.
On every affine scheme chart of the compactification, embed the chart as
a closed subscheme of a smooth affine scheme.  Its closed direct image
of $J_E$ is bounded coherent, and hence compact in the D-module category
of that smooth affine scheme.  By Kashiwara equivalence, an arbitrary
boundary object is a colimit of bounded coherent boundary-supported
D-modules.  Verdier duality and the vanishing of the boundary !-restriction
of the dual of $J_E$ kill maps to each such coherent object.  Compactness
then kills maps to their colimit.  Kashiwara equivalence identifies these
with the mapping complexes on the original, possibly singular chart.
Apply this argument also on scheme charts of every level of a smooth
hypercover.  Smooth descent expresses the stack mapping complex as the
limit of these zero complexes, so it vanishes.  This uses compactness only
in smooth affine ambient charts; no global compactness of $J_E$ is required.

Push by the proper map $\bar p$.  Using $q^*=q^![-2e]$, the simultaneous
extension gives
\[
 p_!q^*E\simeq p_*q^!E[-2e]=L_PE
\]
with its adjunction against every target object.  Compact objects generate
the Levi category.  The functor on the right is continuous, so the same
adjunction extends by colimits to every $E$.  This proves \eqref{eq:M1}.
\end{proof}

\subsection{Second adjointness with the determinant twist}

The first adjunction controls compact Eisenstein objects.  We also need the
opposite-parabolic adjunction, including its unit and counit, to control
objects orthogonal to Eisenstein series.

\begin{proposition}\label{pseudo:second-adjunction}
On the same torsion component there is an adjunction
\begin{equation}\label{eq:M3}
 C_{P^-}\ \dashv\ E_P.
\end{equation}
\end{proposition}

\begin{proof}
Choose an integral cocharacter $\lambda$ defining $P$.  Use the smooth graph
interpolation group $\widetilde H$ over $\mathbb A^1$, with general fiber
ordered as $(g,t^{-\lambda}gt^\lambda)$ and special fiber
$P^-\times_M P$.  The interpolation of bundle stacks is smooth over
$\mathbb A^1$, and its map to $\mathbb A^1\times Y\times Y$ is schematic of
finite type; these are
\cite[Section~2.4 and Proposition~2.5.2]{DrinfeldGaitsgoryCT}.  Let $T$ be the
open substack obtained by retaining degree $u$ at the special fiber and all
of the general fiber.

The determinant-line ratio of the two ends has its canonical splitting for
$t\ne0$.  It extends across $t=0$.  Indeed its order there is locally
constant, and the special fiber is connected: it is a tower of vector
stacks over the connected $Y_M^u$.  Test the order on split reductions,
using the constant diagonal Levi subgroup.  The order is the sum of the
$\lambda$-weights times Euler characteristics, namely $\pm Q(\lambda,u)=0$.
Smoothness of $T$ then extends the section as a nowhere-vanishing section.
Its value at zero equals the parabolic splitting: their quotient is a
unit, is constant along each vector-stack fiber, and equals one on the
split section.  The splitting is equivariant under dilation of the
parameter.  This holds away from zero, where inner conjugation identifies
the induced bundle action with the identity, and hence holds everywhere.

Push $\omega_T$ to $\mathbb A^1\times Y^2$ with this ratio splitting.  The
result is conic in the first factor.  Its !-fiber at one is the identity
kernel, and its !-fiber at zero is the kernel of $E_PC_{P^-}$.  The
cospecialization map between these fibers supplies
\[
 \eta:\operatorname{Id}\longrightarrow E_PC_{P^-}.
\]
To specify the cospecialization and its continuity, on the conic axis use
\[
 i_{1,*}i_1^!\longrightarrow\operatorname{Id}
 \longleftarrow i_{0,*}i_0^!
\]
and apply compactly supported image along $\mathbb A^1$.  Contraction makes
the second arrow invertible.  One can check this on holonomic conic
objects by dualizing: the star extension from the punctured axis has zero
compactly supported image.  Pullback of a compact object from
$[\mathbb A^1/\mathbb G_m]$ is holonomic, because its characteristic variety
is the zero section off zero and has dimension at most one over zero.
Continuous extension from these objects defines the construction on the
conic category used here.  It also tensors with any compactly generated
spectator category, by compactness of the holonomic axis objects.  In these
!-fiber conventions cospecialization for a constant family is the identity.

Define the counit
\[
 \varepsilon:C_{P^-}E_P\longrightarrow\operatorname{Id}
\]
by restricting the kernel on $Y_P^u\times_Y Y_{P^-}^u$ to the open where the
reductions are everywhere transverse.  This open is $Y_M^u$; restriction
followed by star image gives the arrow.  The line identification there is
the split-reduction identification.

We verify the two triangular identities rather than infer them from an
untwisted result.  Precompose $\eta$ with $E_P$.  The family of kernels is
then the fiber product $R=Y_P^u\times_YT$, with $Y_P^u$ on the input end.
Include the general fiber and the everywhere-transverse locus of its
special fiber in an open $R^o$.  It is $\mathbb A^1\times Y_P^u$: an extra
reduction is a flag in $H/P$; its open orbit under the first projection of
$\widetilde H$ has stabilizer identified by the second projection with the
constant $P$.  At zero this is the big orbit of $P^-$, and the graph of
contracting conjugation on $P$ gives the identification across zero.
The maps to the output $H$-bundle and input Levi bundle are constant on
this product.  So is their determinant splitting, since the two
trivializations agree at zero and a unit is constant along an affine line.
Restriction to $R^o$ identifies the unit followed by the counit with
cospecialization of this constant family, which is the identity.

For the other triangular identity use
$T\times_Y Y_{P^-}^u$, with the parabolic factor on the output end.  The
second projection of $\widetilde H$ acts on the flags; on the corresponding
open orbit its stabilizer is identified by the first projection with
$P^-$.  The open family is $\mathbb A^1\times Y_{P^-}^u$, with constant maps
and constant line splitting.  The same calculation gives the identity.
All kernel operations here are safe, or are the holonomic partial
operations described above; their extension with spectators and through
open exhaustions is therefore legitimate.  This proves \eqref{eq:M3}.
\end{proof}

The two adjunctions interact with the extraordinary diagonal as follows:
\begin{equation}\label{eq:M4}
 C_{P^-}F_H\simeq F_M(L'_P)^\vee,
 \qquad F_HE_{P^-}\simeq L_PF_M.
\end{equation}
For the first identity apply $C_{P^-}$ to the output factor of
$\Delta_!k_Y$.  In the following calculation $p,q$ are the maps for $P$.
The expression $q_!p^*$ is the partial left adjoint of $E_P=p_*q^!$
on locally holonomic arguments, with its adjunction against arbitrary
targets.  Proposition~\ref{pseudo:second-adjunction} provides the genuine
left adjoint $C_{P^-}$; uniqueness of left-adjoint values identifies its
value on this kernel with $q_!p^*$.  Base change now gives
$(p,q)_!k_{Y_P^u}$.  Applying $L'_P$ to the first factor of
$\Delta_!k_{Y_M^u}$ gives the same kernel, by \eqref{eq:M1}.
The second identity is its transpose at the opposite twist.  Notice that
$(L'_P)^\vee$ is the right adjoint of $E_P$.  Thus \eqref{eq:M4} identifies
actual continuous functors, not just their values on holonomic test objects.

\subsection{The boundary of the diagonal}

The preceding identities deal with the Eisenstein part of the category.
To complete the proof, we must show that the difference between $F_H$ and
a shift of the identity is built from that part.

\begin{proposition}\label{pseudo:boundary-factorization}
For semisimple $H$, there is a natural transformation
\[
 F_H\longrightarrow\operatorname{Id}[-2\dim Y]
\]
whose cone has a finite filtration with quotients factoring through constant
term, a functor between Levi categories, and Eisenstein series for proper
parabolics.
\end{proposition}

\begin{proof}
We first describe the required compactification, including its properness
for the given global form.  The Vinberg semigroup $\operatorname{Vin}_H$
is the affine Rees monoid of the enhanced group.  Its parameter space is
$T_{\mathrm{ad}}^+\simeq\mathbb A^{\operatorname{rk}H}$, with coordinates
indexed by simple roots.  It has matrices $h_\mu$ in every representation
of dominant highest weight $\mu$ in the \emph{actual} weight lattice of
$H$.  Finitely many generators suffice.  Cartan multiplication has no
parameter factor, while the component of highest weight $\nu$ in the
product for $\mu,\mu'$ is multiplied by $v^{\mu+\mu'-\nu}$.
The nondegenerate locus, where all the $h_\mu$ are nonzero, is smooth over
the parameter space; each fiber is homogeneous for $H\times H$.
At the standard point $c_P$, with coordinates one on Levi roots and zero
on the other simple roots, the matrices are projections to
$V_\mu^{N_P}$, and their stabilizer is $P^-\times_M P$.
These forms of the Vinberg statements, including the smooth stabilizer
family and standard section, are recalled in
\cite[Appendices~C--D]{DrinfeldGaitsgoryCT}.

Let $\bar Y$ classify two $H$-bundles and a section of their associated
Vinberg bundle which is generically nondegenerate on $X$, modulo the
Cartan torus $T_H$.  The parameter $v$ is constant on the complete curve.
Write $b:\bar Y\to Y^2$ for the projection.  Over $v\in T_{\mathrm{ad}}$
there is an open immersion $j_0$ and a finite \'{e}tale $Z_H$-torsor
$l:Y\to\operatorname{im}(j_0)$: fixing $v=1$ gives an isomorphism between
the two bundles, with the residual central scaling.  Thus $bj_0l=\Delta$.

The map $b$ is representable and proper.  Relative finite type follows by
writing the finitely many regular matrix maps $h_\mu$ and their equations.
The torus acts freely relative to the two fixed bundles: all matrices are
generically nonzero, and the dominant weights generate the character
lattice.  For the valuative criterion, fix extended bundles over a DVR.
A torus torsor over its fraction field is trivial, so choose matrix data
there.  The minimum valuation $a(\mu)$ of $h_\mu$ at the generic point of
the special curve is additive under Cartan multiplication.  Indeed the
Cartan product of two nonzero residue matrices is nonzero.  Explicitly,
let $A^+=\bigoplus_\mu V_\mu$ and
$A^-=\bigoplus_\mu V_\mu^*$ be the two Cartan algebras, realized as
coordinate rings of basic affine spaces for opposite Borel conventions.
Inside their tensor product, the diagonal graded subalgebra is
\[
 A=\bigoplus_\mu(V_\mu\otimes V_\mu^*)
   =\bigoplus_\mu\operatorname{End}(V_\mu).
\]
Its product on two matrices $a_\mu,a_{\mu'}$ is
$\pi_{\mu,\mu'}(a_\mu\otimes a_{\mu'})\iota_{\mu,\mu'}$, where
$\iota$ and $\pi$ are the normalized inclusion and projection of the
Cartan component.  This is exactly the matrix product in the Vinberg
Cartan relation.  The basic affine spaces are geometrically integral;
hence the coordinate ring of their product and the displayed subalgebra
are domains, also after
extending to the residue function field.  The nonvanishing assertion
follows.  Hence $a$ extends to
an integral homomorphism from the actual weight lattice to $\mathbb Z$,
that is, a cocharacter of $T_H$.  Rescale by its value on the uniformizer.
Every matrix then has valuation zero and extends over the curve family;
normality removes the remaining apparent poles of these bundle maps.
The parameter coordinates are integral as well.  For a simple root
$\alpha_i$, choose sufficiently large actual dominant weights $\mu,\mu'$
so that the tensor product contains the component
$\mu+\mu'-\alpha_i$.  In its product identity the left side is integral
and the matrix of that component has valuation zero, forcing $v_i$ to be
integral.  The extended matrices are generically nonzero on the special
curve.  The rescaling is unique modulo $T_H$ over the DVR, which proves
uniqueness in the valuative criterion and hence properness.

Stratify $\bar Y$ first by the parameter orbits, indexed by the finitely
many standard parabolics for a fixed Borel.
On the orbit of $c_P$ one may set $v=c_P$, retaining the residual $Z_M$
scaling.  Refine by Levi degrees and defect.  The base $B$ of a refined
stratum consists of two Levi bundles and a positive modification between
them, described by the closure $\bar M$ of $M$ in the Levi blocks of the
semigroup, modulo this residual scaling.  The stratum is the pullback of
\begin{equation}\label{eq:M5}
 q^-\times q:Y_{P^-}\times Y_P\longrightarrow Y_M\times Y_M
\end{equation}
over $B$.  For simply connected derived group this is the defect
stratification in \cite[Proposition~3.2.2]{SchiederVinberg}.  The following
description verifies what we need without changing the global form.

At the generic point of $X$, the homogeneous-orbit description produces
the two parabolic flags.  They extend over the curve by properness of the
flag varieties.  Every matrix factors through its Levi block because it
does so generically; the remaining regular maps give exactly a
$\bar M$-modification.  Conversely these data give the Vinberg section.
In families, take actual positive multiples $\mu_i$ of the fundamental
weights for omitted simple roots.  Their Levi blocks are lines.  Stratify
by lengths of the zero divisors of the generically rank-one matrix maps,
and flatten those zero loci.  On each piece the divisors are relative
Cartier divisors.  Divide out their equations: the resulting image lines
and dual image lines recover the flags.  Their ranks and Pl\"ucker
relations, and all block factorizations, follow from the generically
nondegenerate locus, which is schematically dense on the relative smooth
curve.  Passing to reduced strata suffices for D-modules.  These finitely
many weights embed the flag varieties; the remaining conditions concern
only the modification.  This constructs the claimed locally finite
stratification with fibers \eqref{eq:M5}.

A nonempty defect divisor gives a nonzero rational difference of Levi
degrees, by pairing with these Levi characters.  Conversely, if all these
divisors vanish, determinants of every Levi block are nonzero.  Indeed
powers of these determinants are Laurent monomials in the line characters,
which span the rational character space; after clearing denominators these
relations hold on the closure.  Thus the modification lies in $M$ and has
zero defect.  This proves both implications used below.

On $\bar Y$ take the kernel
\[
 J^b=(j_0)_!l_*k_Y
\]
with the ratio twist of its two ends and the diagonal splitting on the
cover.  Its !-restriction to a stratum is pulled back from $B$.  We give the
argument, since this assertion is the reason the degree vanishing test can
be applied to this kernel.  Choose $z\in X$ where the section is
nondegenerate.  Trivialize the torus quotient locally and, by an \'{e}tale
torus rescaling, set all nonzero parameter coordinates equal to one.
Frame the open-orbit datum at $z$.  The stabilizer group scheme is smooth
over the parameters.  Laurent directions in
$\mathfrak n_{P^-}\oplus\mathfrak n_P$ at $c_P$ therefore extend locally to
sections of its Lie algebra, using only finitely many jets for a chosen
finite set of directions.  Infinitesimal gluing at $z$ defines derivations
on smooth presentations of the ambient Vinberg bundle stack.  They
preserve $v$, hence the pair defining $j_0$, and lift compatibly to the
stratum fiber over $B$: there they change the two unipotent extensions and
keep the Levi modification fixed.  Principal parts on the two factors,
with the relative chart directions, span that fiber's tangent space.
The finite \'{e}tale image $l_*k_Y$ is a local system on the open.  Apply
the analytic-flow observation and then vector-stack descent exactly as in
Lemma~\ref{pseudo:fiber-descent}.  This proves the assertion for $J^b$.

On $B$ there are independent central-inertia cocharacters at its two ends:
the residual $Z_M$ quotient absorbs the difference of the two central
actions on the modification.  Their weights on the ratio line are
$Q(z,u_1)$ and $-Q(z,u_2)$.  The irrational Kummer test therefore kills the
twisted category unless both $u_1$ and $u_2$ are torsion.  Their difference
then vanishes rationally, so the defect must be zero.  Only the defect-zero
strata of torsion degree remain.  There are finitely many such degrees and
finitely many parabolics.

Set
\[
 D^b=\operatorname{cone}\bigl(J^b\longrightarrow(j_0)_*l_*k_Y\bigr).
\]
For a boundary stratum inclusion $i$, the identity
$i^!(j_0)_*=0$ gives $i^!D^b\simeq i^!J^b[1]$.  Thus the preceding
boundary calculation applies to the cone.  Use localization with
!-restrictions and star extensions from strata to filter $D^b$.
Within a fixed parameter orbit, the !-restrictions on every nonzero-defect
stratum vanish by the preceding argument; this is checked on
quasi-compact pieces using the locally finite stratification.  Thus the
cone has a finite filtration involving only the defect-zero strata for
proper parabolics.  Push it by the proper map $b$.  On a stratum over $B$
its contribution has the form
\[
 (p^-\times p)_*(q^-\times q)^!K
\]
for a kernel $K$ on $Y_M\times Y_M$, by !-base change from $B$.  The map
forgetting the modification is quasi-compact and safe on fixed degrees:
in the remaining defect-zero case it is the isomorphism datum modulo
central scaling, which acts freely relative to the two Levi bundles.
If $T_K$ denotes the functor represented by $K$, the displayed kernel
represents $E_P\circ T_K\circ C_{P^-}$; any shift from the cone is
absorbed into $K$.  This is the required factorization through both a
proper constant term and a proper Eisenstein image.  Finally, $bj_0l=\Delta$, $b$ is proper, and
$k_Y=\omega_Y[-2\dim Y]$.  The two pushed kernels are therefore
$\Delta_!k_Y$ and $\Delta_*\omega_Y[-2\dim Y]$, respectively.  This proves
the Proposition.
\end{proof}

\subsection{Induction and the reductive case}

We finish the proof of Theorem~\ref{pseudo:equivalence}.  The induction is
on semisimple rank, simultaneously for the two signs of $s$.  We first
explain the semisimple step, assuming the assertion for reductive groups
of smaller semisimple rank.

For every proper parabolic and torsion degree, choose compact generators
$w$ of $\mathcal C_M^u$ and put $x=E_Pw$.  Equation~\eqref{eq:M1} implies
that $x$ is compact.  Equation~\eqref{eq:M4} gives
$F_Hx\simeq L_{P^-}F_Mw$, which is compact by the inductive equivalence.
For every $y$, adjunction and \eqref{eq:M4} give natural isomorphisms
\[
 \operatorname{Hom}(F_Hx,F_Hy)
 \simeq \operatorname{Hom}(F_Mw,C_{P^-}F_Hy)
 \simeq \operatorname{Hom}(w,(L'_P)^\vee y)
 \simeq \operatorname{Hom}(x,y).
\]
Full faithfulness requires the actual map induced by $F_H$ to be
invertible.  The preceding adjunction calculation gives an isomorphism
natural in $y$, but has not yet shown that its inverse sends the identity
of $x$ to the identity of $F_Hx$.  By enriched
Yoneda, the inverse of the displayed natural isomorphism has the form
$f\mapsto F_H(f)b$ for an endomorphism $b$ of $F_Hx$.  Surjectivity at $y=x$
provides $a\in\operatorname{End}(x)$ with $F_H(a)b=1$.  A split epimorphism
endomorphism of $F_Hx$ is invertible: on affine charts its bounded coherent
cohomology modules are noetherian, and a surjective endomorphism of a
noetherian module is injective.  Hence $F_H(a)$ and then $b$ are invertible.
The actual $F_H$-comparison is consequently an isomorphism.

Let $\mathcal B$ be the right orthogonal of this set of compact $x$.
Equation~\eqref{eq:M1} identifies it with the objects whose proper
constant terms vanish.  Equation~\eqref{eq:M3} also makes these objects
left orthogonal to every proper Eisenstein image.  If $m\in\mathcal B$,
Proposition~\ref{pseudo:boundary-factorization} therefore gives both
$F_Hm\simeq m[-2\dim Y]$ and
\[
 \operatorname{Hom}(m,y)\xrightarrow{\ \sim\ }
 \operatorname{Hom}(F_Hm,F_Hy).
\]
The second assertion follows by applying $\operatorname{Hom}(m,-)$ to
the cone filtration and using naturality of its map.  Compact localization
says that the $x$ together with $\mathcal B$ generate under colimits.
Mapping out of colimits converts them to limits, so these comparisons
prove full faithfulness of $F_H$ on every source.  Its essential image is
closed under colimits.  It contains $\mathcal B$ because $F_H$ is a shift
there, and it contains every Eisenstein generator by \eqref{eq:M4} and
the inductive equivalence.  Hence it is the entire category.

It remains to justify the reductive groups used at the induction base and
as Levi subgroups.  For a reductive $H$, let $H_{\rm der}$ be its derived
group and $Z^0$ its connected center.  Their product maps to $H$ by a
central isogeny.  On torsion degrees, the induced map of bundle stacks
\[
 f:Y_{H_{\rm der}}\times Y_{Z^0}^{\,0}\longrightarrow Y_H^{\rm tors}
\]
retaining all contributing lifting components is surjective.  In fact the
abelianization degree of a torsion class is zero.  In the exact sequence
\[
 \pi_1(H_{\rm der})\longrightarrow\pi_1(H)
 \longrightarrow\pi_1(H/H_{\rm der})\longrightarrow0,
\]
the left image is finite and the right group is a free lattice, so the
left image is exactly the torsion subgroup.  Thus the degree lifts from
the derived group.  The degree-two obstruction for the finite central
isogeny is the class of this degree in the cokernel of the map on
fundamental groups; it therefore vanishes.  One may check this last statement topologically on the complex
curve and then use comparison for finite coefficients.  When a lift
exists, its choices form a torsor for bundles of this finite kernel.
There are finitely many such bundles and finite automorphism groups.
The map $f$ is therefore, \'{e}tale locally on the base, a finite union of
finite \'{e}tale gerbes.  Formal \'{e}taleness follows as well from the
isomorphism of Lie algebras.  In characteristic zero, $f$ is proper and
safe, $f^!=f^*$ is conservative, and pushforward and pullback are
ambiadjoint.

The twisting separates on this product.  Decompose $V$ into representations
of $H_{\rm der}$ tensored with central characters.  The determinant-of-
cohomology tensor formula separates each summand into its derived and
central factors; its mixed Deligne-pairing term vanishes because a
representation of a semisimple group has trivial determinant.  On the
neutral component of the torus stack, a rigidification at a point of $X$
identifies the stack with a Picard variety times $BZ^0$.  Writing $V=\bigoplus_i W_i\otimes\chi_i$, its central inertia weight
on degree zero is $(1-g(X))\sum_i\dim(W_i)\chi_i=0$, because $V$ is
self-dual.  Thus only the Picard factor is twisted.  Its
pseudo-identity is a shift, since the variety is smooth and separated.
For $BZ^0$ use the untwisted miraculous duality theorem
\cite[Theorem~0.1.6]{GaitsgoryStrange}, applied to the neutral component of
$\operatorname{Bun}_{Z^0}^{\,0}(\mathbb P^1)=BZ^0$.  The extraordinary diagonal kernel of the product is the exterior
product of those of its factors, so its pseudo-identity is
$F_{H_{\rm der}}\otimes F_{Z^0}$.  Thus rank zero is established, and
the semisimple assertion of any rank gives the product assertion of that
rank.

Finally that assertion descends along $f$.  The extraordinary diagonal
kernels give
\[
 f_*F_{H_{\rm der}\times Z^0}\simeq F_Hf_* ,\qquad
 F_{H_{\rm der}\times Z^0}f^!\simeq f^!F_H.
\]
For example both kernels for the first equality are
$\Gamma_{f,!}k$, where $\Gamma_f$ is the graph: use $f_*=f_!$ and then
star--shriek base change with $f^!=f^*$, which holds because $f$ is
\'{e}tale.  Transposition at the opposite twist gives the second equality.
Compact objects pushed from the product remain compact and generate
downstairs, since $f^!$ is continuous and conservative.  Their images
under $F_H$ are compact by the product assertion.  The displayed identities
and ambiadjunction give the same natural Hom isomorphism as above on these
generators.  The noetherian split-epimorphism argument identifies it with
the actual comparison map.  Full faithfulness follows by generation;
essential surjectivity follows because the pushed generators lie in the
image.  This completes the simultaneous induction and proves
Theorem~\ref{pseudo:equivalence}.
